% GENERATED FILE. Edit canonical lessons, then run npm run generate:books.
% canonical-source-hash: fnv1a64:beb32b8b66098647
% canonical-lessons: ML-C08-dayavayi

\chapter{Please}
\label{ch:ml-please}

This chapter is generated from the canonical micro-lessons used by Language Ladder.
Each section stays independently resumable and preserves the authored prerequisite order.

\section[\ml{ദയവായി}]{\ml{ദയവായി} (dayavāyi) — please (as a compassion)}

\label{lesson:ML-C08-dayavayi}



\begin{quote}
\textbf{Warm-up.} {[}PAUSE 2s{]} Malayalam's "please" turns compassion into an adverb: do this thing \textbf{as a kindness}.
\end{quote}

\subsection*{You'll want to know}
\textbf{\ml{ദയവായി}} = \textbf{please}, literally "\textbf{kindly / as a compassion}." Two pieces:

\begin{itemize}
\raggedright
  \item \textbf{\ml{ദയ}} (\emph{daya}) = \textbf{compassion, kindness, grace} — from Sanskrit \textbf{daya}.
  \item \textbf{‑\ml{വായി}} (\emph{-vāyi}) = "\textbf{having become / as}," turning the noun into "\emph{in the manner of} compassion" — i.e. \textbf{kindly}. (The true historical seam is \emph{dayavu} "kindness" + \textbf{‑\ml{ആയി}} \emph{-āyi}, the adverb of \textbf{\ml{ആകുക}} \emph{ākuka} "to become"; we split it \emph{-vāyi} here just to keep the \emph{daya} root visible.)
\end{itemize}

So \emph{dayavāyi} asks you to act \emph{kindly}, \emph{as a kindness}. Where Tamil and Telugu say "\textbf{do} the kindness" and Kannada "having \textbf{placed} it," Malayalam makes kindness the \textbf{manner} of the whole request.

The family resemblance is unmistakable, and it reaches beyond Dravidian: \textbf{Arabic} \emph{min faḍlik} ("from your \textbf{grace}") and \textbf{Hindi} \emph{kṛpayā} (from \emph{kṛpā}, "\textbf{compassion}") ask by the same grace-naming instinct.

\begin{culture}
\textbf{daya} + a light verb, four ways:

\begin{itemize}
\raggedright
  \item Malayalam \textbf{\ml{ദയവായി}} \emph{daya\textbf{vāyi}} — compassion \textbf{+ as / having become}
  \item Tamil \textbf{\ta{தயவுசெய்து}} \emph{tayavu \textbf{sey}tu} — kindness \textbf{+ do}
  \item Kannada \textbf{\ka{ದಯವಿಟ್ಟು}} \emph{daya\textbf{viṭṭu}} — compassion \textbf{+ having placed}
  \item Telugu \textbf{\te{దయచేసి}} \emph{daya\textbf{cēsi}} — compassion \textbf{+ having done}
\end{itemize}
\end{culture}

\begin{culture}
Honestly: as elsewhere in the family, the standalone word is a touch \textbf{formal / emphatic}. Everyday Malayalam politeness leans on the \textbf{respectful request} form of the verb — endings like \textbf{‑\ml{ൂ}} (\emph{-ū}) or \textbf{‑\ml{ണം}} (\emph{-aṇaṁ}): \textbf{\ml{ഇരിക്കൂ}} (\emph{irikkū}) "please sit," \textbf{\ml{പറയൂ}} (\emph{paṟayū}) "please tell me." Add \textbf{\ml{ദയവായി}} in front when you want the courtesy spoken aloud.
\end{culture}

\begin{cousinweb}
The heart of the word is \textbf{\ml{വാ}} — \emph{va} stretched by the \textbf{long‑ā} sign \textbf{\ml{ാ}} (\emph{-ā}) — then \textbf{\ml{യി}}, \emph{ya} + the \emph{i} sign (\textbf{\ml{ി}}). Malayalam's rounded letters hang their vowel-signs to the right; the long \textbf{\ml{ാ}} is simply a tall stroke after the consonant, the same \emph{ā} you would meet across the Dravidian scripts.
\end{cousinweb}

\subsection*{Guided Practice}
{[}PAUSE 1s{]}

\begin{itemize}
\raggedright
  \item {[}YOU SAY: "daya" — compassion — then "-vāyi," as / kindly{]}
  \item {[}YOU SAY: the whole word — "dayavāyi" = please{]}
  \item {[}YOU SAY: the everyday way — the verb ending: "irikkū" = please sit{]}
\end{itemize}

\begin{tcolorbox}[breakable,colback=teal!4,colframe=teal!35!black,title={Wrap-up recall}]
{[}PAUSE 3s{]} What is the Malayalam word for please? (\textbf{\ml{ദയവായി}} \emph{dayavāyi}.) What do its pieces mean? ("\textbf{Compassion}" \emph{daya} + "\textbf{as / having become}" \emph{-vāyi} — \emph{kindly}.) Which cousins ask the same way? (\textbf{Tamil, Kannada, Telugu} with \emph{daya} — plus \textbf{Arabic} \emph{faḍl}, \textbf{Hindi} \emph{kṛpā}.) What carries "please" in ordinary speech? (\textbf{The respectful verb form}, endings like \textbf{‑\ml{ൂ}} \emph{-ū} / \textbf{‑\ml{ണം}} \emph{-aṇaṁ}.)
\end{tcolorbox}
